\problemname{Robber's Language}
In the Robber's language, text is translated by replacing every consonant $x$ with $x$\texttt{o}$x$, while all vowels are left unchanged.
For example, \texttt{hej} is translated to \texttt{hohejoj}, \texttt{moo} to \texttt{momoo}, and \texttt{ojoj} to \texttt{ojojojoj}.

Oskar has just learned the Robber's language, but sometimes forgets to transform certain consonants. Oskar thus only transforms a \emph{subset} of the consonants in the manner described above.
You will be given an original text and a potential translation that Oskar has made. You should determine whether it is possible that Oskar translated the first string into the second according to the description above.

If Oskar had written \texttt{hejoj}, the original text could for example have been \texttt{hej} or \texttt{hejoj}.
For clarity, we define the following letters as vowels: \texttt{a,e,i,o,u,y}, and all other letters are consonants.

\section*{Input}
The first line contains the word that we wonder if it could have been the original text. The second line contains a word that Oskar has written. Both words contain at least one character and only characters \texttt{a-z} in lowercase. No string is longer than 1000 characters.

\section*{Output}
The output should consist of a single word: \texttt{ja} if it is possible that the second word is a translation of the first, otherwise \texttt{nej}.
The answer should thus be \texttt{ja} if and only if it is possible to choose some subset of the consonants in the first word and apply the Robber's language transform to get the second word.

\section*{Scoring}
Your solution will be tested on a set of test groups.
To earn points for a group, you must pass all test cases in that group.

\noindent
\begin{tabular}{| l | l | p{12cm} |}
  \hline
  \textbf{Group} & \textbf{Points} & \textbf{Constraints} \\ \hline
  $1$    & $20$        & The strings are at most 20 characters long \\ \hline 
  $2$    & $30$        & The strings are at most 100 characters long \\ \hline
  $3$    & $50$        & No additional constraints. \\ \hline 
\end{tabular}

\section*{Explanation of Samples}
Explanation of sample 1: it does not matter which consonant Oskar chooses to apply the Robber's language transform to --- as long as
he chooses exactly one of the \texttt{j}'s. Say Oskar forgot to transform all \texttt{j}'s except the first --- then \texttt{ojojoj} becomes \texttt{ojojojoj}. The answer is thus ``\texttt{ja}''.

Explanation of sample 2: it is impossible to transform the first string into the second. One way to show this is to observe that in the first string there are two \texttt{j}'s next to each other, which means there will always be two \texttt{j}'s next to each other. Thus it is not possible to transform the first string into the second, and the answer is ``\texttt{nej}''.
